\title{FlashAttention-3: Fast and Accurate Attention with Asynchrony and Low-precision}

\begin{abstract}
Attention forms a foundational element in the widely adopted Transformer framework, but it also poses a significant computational bottleneck for scaling large language models and extending context lengths. \textsc{FlashAttention} introduced an efficient GPU-based approach that reduces the overhead associated with memory access during attention computation. Yet, this design does not fully leverage the latest hardware advancements; for example, \textsc{FlashAttention-2} utilizes just 35\% of the H100 GPU's peak capacity. To address this, we introduce three core innovations tailored for Hopper GPUs: (1) utilizing the asynchrony between Tensor Cores and TMA to both synchronize data transfer and computation through warp-specialization, (2) interleaving block-level matrix multiplications and softmax operations, and (3) applying block quantization and non-uniform processing strategies that tap into dedicated hardware support for FP8 low-precision operations. Our solution, \textsc{FlashAttention-3}, attains 1.5-2.0$\times$ greater throughput on H100 GPUs, achieving up to 740 TFLOPs/s in FP16 (representing 75\% utilization), and approaching 1.2 PFLOPs/s with FP8 arithmetic. Furthermore, we provide evidence that FP8 \textsc{FlashAttention-3} produces a numerical error that is 2.6$\times$ less than that observed in standard FP8 attention baselines.
\end{abstract}

\section{Introduction}
\label{sec:intro}

Within the Transformer architecture~\citep{vaswani2017attention}, the computation of self-attention, particularly the calculation of attention scores between queries and keys, represents the principal performance constraint due to its quadratic dependence on sequence length. Extending the reach of attention mechanisms to handle much longer contexts could enable a range of new functionalities—including improved reasoning across numerous lengthy documents~\citep{guo2021longt5,shaham2022scrolls,peng2023yarn}, more effective navigation and understanding of extensive codebases~\citep{roziere2023code, li2023starcoder}, as well as accommodating new data modalities like high-resolution images~\citep{chen2022scaling}, audio streams~\citep{gulati2020conformer}, and video~\citep{ho2022video}. Additionally, this extension may benefit applications involving prolonged user histories~\citep{sun2019bert4rec} and long-horizon agent workflows~\citep{yao2022react}. These possibilities have motivated considerable research on accelerating attention for long sequences, utilizing approaches such as approximation techniques~\citep{katharopoulos2020transformers,choromanski2020rethinking,
tay2020efficient}, advancements in software efficiency~\citep{rabe2021self,dao2022flashattention,kwon2023efficient}, and the exploration of alternative neural architectures~\citep{peng2023rwkv,sun2023retentive,gu2023mamba}.

In this paper, we extend the efforts of \citet{dao2022flashattention}, who devised exact-attention algorithms that intentionally leverage an understanding of the GPU’s architectural details and execution paradigm to inform algorithmic structure. The seminal work in \citep{dao2022flashattention} presented \textsc{FlashAttention}, introducing an innovative tiling approach that consolidates all components of the attention mechanism into a unified GPU kernel, thus eliminating the overhead associated with multiple intermediate accesses to slow global memory.

Subsequently, \citet{dao2023flashattention2} reformulated the algorithm as \textsc{FlashAttention-2}, enhancing parallelism by operating across the sequence length dimension and structuring the core of the forward pass around chunks of the key and value matrices. This reorganization aimed to achieve more effective workload allocation and resource use on the GPU. Nonetheless, we note that \textsc{FlashAttention-2} underutilizes computational resources on more recent GPU architectures when compared to highly optimized matrix multiplication (GEMM) kernels—demonstrating, for instance, only 35\% utilization versus 80–90\% on the Hopper H100 GPU. A contributing factor to this disparity is the persistence of older, Ampere-optimized instructions instead of Hopper-native alternatives for operations targeting Tensor Cores. Other recent work, such as ThunkerKitten~\citep{spector2024thunder} and cuDNN 9~\citep{cudnn9}, has illustrated that incorporating Hopper-specific instructions and employing tile-based representations can both accelerate attention operations and streamline their implementation.

At a deeper level, the algorithm underlying \textsc{FlashAttention-2} is based on a straightforward, synchronous computation model and does not incorporate asynchronous execution or low-precision arithmetic within its architecture. Hardware asynchrony arises from the design of specialized processing components tailored to enhance critical machine learning operations: certain units handle matrix multiplications (such as Tensor Cores) or memory operations (like the Tensor Memory Accelerator—TMA), functioning independently from general-purpose CUDA cores, which execute logical, integer, and floating-point instructions. The adoption of lower precision formats—including FP8 on Hopper and FP4 on Blackwell, building on previous transitions to FP16 (Pascal, 2017) and BF16 (Ampere, 2020)—remains an effective means of doubling or quadrupling computational throughput without increasing energy or silicon footprint. The enhancements provided by Hopper in these respects are examined in greater depth in \cref{subsec:hardware}. The main technical obstacle is reformulating \textsc{FlashAttention-2} to leverage these advanced hardware functionalities: implementing asynchrony demands carefully orchestrating overlapping execution of matrix multiplications and softmax—despite their data dependencies—while supporting low-precision computation necessitates strategies for mitigating quantization errors, particularly when addressing extreme-valued features in LLMs~\citep{dettmers2208llm, sun2024massive}.

For this purpose, we introduce \textsc{FlashAttention-3}, which integrates and advances three innovative concepts aimed at enhancing efficiency on contemporary GPU platforms.\footnote{Our experiments are conducted on NVIDIA's Hopper architecture for demonstration, but the approach applies broadly to any GPU system with strong support for asynchronous processing and low-precision operations.}

\begin{enumerate}

\item \textbf{Producer-Consumer asynchrony:} We propose a warp-specialized software pipelining approach that leverages the decoupled execution of data movement operations and Tensor Core computations. By assigning producer and consumer roles to separate warps, this design enables more effective latency hiding for both memory accesses and instruction issue stalls in the algorithm.
\item \textbf{Hiding softmax under asynchronous block-wise GEMMs:} We coordinate the execution of the relatively low-throughput non-GEMM computations required for softmax—including floating point multiply-add and exponential functions—with asynchronous WGMMA instructions responsible for GEMM. This entails revising the \textsc{FlashAttention-2} procedure to eliminate certain serial dependencies between softmax and GEMM operations. For instance, in the two-stage variant of our method, while the softmax step is processed for one segment of the scores matrix, WGMMA asynchronously computes results for the subsequent segment.
\item \textbf{Hardware-accelerated low-precision GEMM:} We modify the forward pass routine to utilize FP8 Tensor Cores in GEMM operations, leading to a nearly twofold increase in observed TFLOPs/s. Addressing this requires resolving the differing layout expectations of WGMMA for blocks in FP32 accumulators and FP8 operand matrices stored in memory. By applying block quantization and incoherent processing, we control the degradation in accuracy that may arise when transitioning to FP8 calculations.

To empirically assess the effectiveness of our approach, we conduct performance benchmarks of \textsc{FlashAttention-3} on the H100 SXM5 GPU across a variety of parameter settings. Our results indicate: (1) FP16 delivers a 1.5-2.0$\times$ acceleration relative to \textsc{FlashAttention-2} in the forward pass—reaching peak throughput of up to 740 TFLOPs/s—and achieves a 1.5-1.75$\times$ speedup during the backward pass; (2) FP8 attains nearly 1.2 PFLOPs/s throughput; and (3) at larger sequence lengths, FP16 surpasses the competing methods, while FP8 performs comparably\footnote{Specifically, with a head dimension of 64, FP8 \textsc{FlashAttention-3} leads, whereas for head dimensions of 128 and 256, performance matches that of competitors without causal masking and trails behind when causal masking is applied.} with the current state-of-the-art attention implementation in NVIDIA's cuDNN library. We further confirm that FP16 \textsc{FlashAttention-3} maintains numerical accuracy equivalent to \textsc{FlashAttention-2}, and surpasses the standard attention operation, owing to storage of intermediate results (such as softmax rescaling) in FP32 precision. Additionally, our FP8 \textsc{FlashAttention-3} configuration—employing block quantization and incoherent processing—achieves a 2.6$\times$ accuracy improvement over the conventional attention baseline using per-tensor quantization in scenarios with outlier features.

We have released \textsc{FlashAttention-3} under a permissive open-source license\footnote{\textsc{FlashAttention-3}
  is available at \texttt{https://github.com/Dao-AILab/flash-attention}}, and intend to incorporate it into both PyTorch and Hugging Face frameworks to maximize accessibility for the research and developer communities.

\section{Background: Multi-Head Attention and GPU Characteristics}
\label{sec:background}

\subsection{Multi-Head Attention}
\label{subsec:multi_head_attn}

Let $\mathbf{Q}, \mathbf{K}, \mathbf{V} \in \mathbb{R}^{N \times d}$ be the query, key and value input sequences associated to a single head, where $N$ is the sequence length and $d$ is the head dimension. Then the attention output $\mathbf{O}$ is computed as:

\begin{equation*}
  \mathbf{S} = \alpha \mathbf{Q} \mathbf{K}^\top \in \mathbb{R}^{N \times N}, \quad \mathbf{P} = \mathrm{softmax}(\mathbf{S}) \in \mathbb{R}^{N \times N}, \quad \mathbf{O} = \mathbf{P}\mathbf{V} \in \mathbb{R}^{N \times d},
\end{equation*}

Here, $\mathrm{softmax}$ operates along each row, and it is common to choose the scaling coefficient as $\alpha = 1/\sqrt{d}$. To avoid numerical issues associated with the exponential computation, one typically subtracts $\mathrm{rowmax}(\mathbf{S})$ from $\mathbf{S}$. In the multi-head attention (MHA) framework, individual heads employ distinct projections for the queries, keys, and values. These computations are performed concurrently over all heads and batches, resulting in the aggregated output tensor.

Now let $\phi$ be a scalar loss function and let $\mathbf{d}(-) = \partial \phi / \partial (-)$ be notation for the gradient. Given the output gradient $\mathbf{dO} \in \mathbb{R}^{N \times d}$, we compute $\mathbf{dQ}$, $\mathbf{dK}$, and $\mathbf{dV}$ according to the chain rule as follows:

\begin{align*}
  \mathbf{dV} &= \mathbf{P}^\top \mathbf{dO} \in \mathbb{R}^{N \times d} \\
  \mathbf{dP} &= \mathbf{dO} \mathbf{V}^\top \in \mathbb{R}^{N \times N} \\
  \mathbf{dS} &= \mathrm{dsoftmax} (\mathbf{dP}) \in \mathbb{R}^{N \times N} \\
  \mathbf{dQ} &= \alpha \mathbf{dS} \mathbf{K} \in \mathbb{R}^{N \times d} \\
  \mathbf{dK} &= \alpha \mathbf{dS}^\top \mathbf{Q} \in \mathbb{R}^{N \times d},
\end{align*}

In this context, $\mathbf{d}s = (\mathrm{diag}(p) - p p^\top)\mathbf{d}p$ holds when $p = \mathrm{softmax}(s)$, considering $s$ as an input vector; we denote the row-wise application of this relation by $\mathrm{dsoftmax}(\mathbf{dP})$. Moreover, this process remains parallelizable over both the batch and head dimensions during the backward phase of MHA.

\subsection{GPU hardware characteristics and execution model}
\label{subsec:hardware}

We outline the elements of the GPU execution model pertinent to \textsc{FlashAttention-3}, specifically examining the NVIDIA Hopper architecture as a representative example of this framework.

\paragraph{Memory hierarchy:} The GPU features a layered organization of memory types, where storage size is inversely proportional to access bandwidth (\cref{tab:gpu-hierarchy})\footnote{\citet{luo2024benchmarking} indicates that shared memory achieves a bandwidth of 128 bytes per clock cycle per SM; multiplying this by 132 SMs and the boost frequency of 1830 MHz yields the final throughput.}. At the largest and slowest tier, global memory (GMEM)—also referred to as HBM—serves as off-chip DRAM that all streaming multiprocessors (SMs) can address. Data from GMEM are automatically buffered within an on-chip L2 cache. Each SM is equipped with a limited amount of shared memory (SMEM): an on-chip, highly banked cache that is directly managed by the programmer. The fastest access is provided by the register file specific to each SM.

\paragraph{Thread hierarchy:} In the GPU programming paradigm, execution units known as threads are structured into a multi-tier hierarchy. This begins at the level of individual threads and aggregates upward into warps (each consisting of 32 threads), followed by warpgroups (which contain 4 adjacent warps), then threadblocks—also referred to as cooperative thread arrays (CTAs)—and threadblock clusters (introduced with Hopper architectures), culminating at the grid level.

The connection between these two hierarchies is significant. Threads belonging to a single CTA are assigned to execute together on the same SM, while CTAs grouped within a cluster are executed concurrently on the same GPC. All threads inside a CTA can directly access SMEM, but each thread is limited to 256 private registers (RMEM) that are inaccessible to others.

\paragraph{Asynchrony and warp-specialization:}

GPUs function as throughput-oriented processors that use parallelism and asynchronous operations to mask delays associated with memory access and computation. In the case of asynchronous data transfers between GMEM and SMEM, Hopper introduces the Tensor Memory Accelerator (TMA), a specialized hardware component \cite[\S7.29]{cuda}. Additionally, distinct from earlier designs like Ampere, Hopper's Tensor Core—which can be utilized via the warpgroup-wide WGMMA instruction \cite[\S9.7.14]{ptx}—also operates asynchronously and is capable of retrieving inputs straight from shared memory.

The presence of hardware mechanisms for asynchrony enables the implementation of warp-specialized kernels, in which the warps belonging to a single CTA are assigned dedicated roles as either producers or consumers—engaged solely in data transfer or computation, respectively. This configuration generally enhances the compiler’s ability to optimize instruction scheduling \citep{warp-specialization-2011}. Furthermore, Hopper introduces the capability for registers to be dynamically redistributed across warpgroups using \verb|setmaxnreg| \cite[\S9.7.17.1]{ptx}. As a result, warps performing matrix multiply-accumulate (MMA) operations can be allotted a greater portion of RMEM, while warps limited to performing TMA require only minimal resources—potentially as little as one participating thread.

\paragraph{Low-precision number formats:}
\label{sec:low-precision-gpu}
Contemporary GPUs are equipped with dedicated hardware specifically designed to boost the efficiency of low-precision operations. As an illustration, the WGMMA instruction enables the utilization of FP8 Tensor Cores on the Hopper architecture, thereby achieving double the throughput per SM relative to the FP16 or BF16 formats.

Nevertheless, employing FP8 WGMMA properly requires attention to the memory layout restrictions imposed on its inputs. Consider a GEMM operation computing $A \times B^{\top}$, where $A$ is an $M\times K$ matrix and $B$ is $N\times K$. We define an operand as \emph{mn-major} if the elements are contiguous with respect to the outer $M$ or $N$ dimension, while it is considered \emph{k-major} when contiguity is instead along the inner $K$-dimension. For FP16 WGMMA, SMEM operands are permissible in either mn-major or k-major layouts. In contrast, FP8 WGMMA is limited to accepting only k-major operands. Furthermore, scenarios such as attention mechanisms, which benefit from chaining multiple GEMMs within one kernel, are complicated by incompatible layouts between FP32 accumulators and FP8 operands, thereby making it challenging to chain FP8 WGMMA operations with dependencies.

Within the framework of attention mechanisms, these layout constraints necessitate specific adjustments to how an FP8 algorithm is constructed, as detailed in \cref{sec:algofp8}.

\subsection{Standard Attention and Flash Attention} As described by \citet{dao2022flashattention}, \textbf{standard attention} refers to the approach in which attention computation on the GPU involves writing the intermediate matrices $\mathbf{S}$ and $\mathbf{P}$ to HBM. The principal innovation of \textsc{FlashAttention} was to employ a local softmax reduction, thereby minimizing costly intermediate memory transactions and enabling the entire attention computation to be handled by a single fused kernel. In this context, the local softmax aligns with lines \ref{code-ws:softmax_start}-\ref{code-ws:softmax_end} from the consumer mainloop presented in \cref{alg:flash3_wgmma_ws_only}, along with the associated block-wise rescaling operations on $\mathbf{O}$. A concise proof demonstrating that this approach correctly yields $\mathbf{O}$ can be found in \cite[\S 2.3.1]{dao2023flashattention2}.

\section{FlashAttention-3: Algorithm}
\label{sec:algo}

In this section, we present an overview of the \textsc{FlashAttention-3} algorithm. To maintain clarity, our discussion is restricted to the forward pass, while details regarding the backward pass are provided in~\cref{sec:algo_ws_bwd}. We begin by outlining how warp-specialization and a circular SMEM buffer can be incorporated into the foundational procedure of \textsc{FlashAttention-2}. Next, we discuss the use of WGMMA asynchrony to implement a two-stage GEMM-softmax pipeline with overlapping computation. Finally, we detail the changes required to support FP8, considering both layout compatibility and enhancements to precision through block quantization and methods for dealing with incoherent computation.

\subsection{Producer-Consumer asynchrony through warp-specialization and pingpong scheduling}
\label{sec:algo_ws}

\paragraph{Warp-specialization}
Similar to the approach used in \textsc{FlashAttention-2}, the forward computation in \textsc{FlashAttention-3} is highly parallelizable across the dimensions of batch size, head count, and query sequence length. As a result, it is sufficient to present an algorithmic overview at the CTA granularity, wherein the algorithm processes a given tile $\mathbf{Q}_i$ from the query matrix and computes its corresponding output tile $\mathbf{O}_i$. For clarity, we initially describe the warp-specialization method that employs a circular SMEM buffer but does \emph{not} implement GEMM-softmax overlap. Denoting the head dimension by $d$ and the sequence length by $N$, and fixing a block size $B_r$ for queries, we partition $\mathbf{Q}$ into $T_r = \lceil \frac{N}{B_r} \rceil$ distinct blocks, labeled $\mathbf{Q}_1, \ldots, \mathbf{Q}_{T_r}$.

\begin{algorithm}[H]
    \caption{\small\label{alg:flash3_wgmma_ws_only}\textsc{FlashAttention-3} forward pass \textbf{without} intra-consumer overlapping -- CTA perspective}
    \begin{algorithmic}[1]
\REQUIRE Matrices $\mathbf{Q}_i \in \mathbb{R}^{B_r \times d}$ and $\mathbf{K}, \mathbf{V} \in \mathbb{R}^{N \times d}$ reside in HBM, with a specified key block size $B_c$ such that $T_c = \lceil \frac{N}{B_c} \rceil$.
\STATE Create a pipeline handler to coordinate barrier synchronization using a circular SMEM buffer of $s$ stages.
\IF {warpgroup is a producer}
\STATE Free up a specified quantity of registers.
\STATE Begin data transfer of $\mathbf{Q}_i$ from HBM to shared memory.
\STATE After load completes, commit and signal consumer warpgroups of the availability of $\mathbf{Q}_i$.
\FOR{$0 \le j < T_c$}
    \STATE Wait until the $(j\,\%\,s)$ stage in the buffer has been consumed.
    \STATE Start loading $\mathbf{K}_j, \mathbf{V}_j$ from HBM into shared memory at the $(j\,\%\,s)$ buffer location.
    \STATE On load completion, commit and inform consumers that $\mathbf{K}_j, \mathbf{V}_j$ are available.
\ENDFOR
\ELSE \STATE Allocate a defined number of registers, determined by how many consumer warps are present.
\STATE Initialize on-chip: $\mathbf{O}_i = (0) \in \mathbb{R}^{B_r \times d}$ and $\ell_i, m_i = (0), (-\infty) \in \mathbb{R}^{B_r}$.
\STATE Wait for $\mathbf{Q}_i$ to arrive in shared memory.
\FOR{$0 \le j < T_c$}
\STATE Wait until $\mathbf{K}_j$ becomes available in shared memory.
\STATE Calculate $\mathbf{S}_i^{(j)} = \mathbf{Q}_i \mathbf{K}_j^T$ using SS-GEMM. Commit, then block until operation finishes. \label{alg:ws_only_gemm1}
\STATE Save $m_i^{\mathrm{old}} = m_i$ and update $m_i = \mathrm{max}(m_i^{\mathrm{old}}, \mathrm{rowmax}(\mathbf{S}_i^{(j)}))$. \label{code-ws:softmax_start}
\STATE Compute $\widetilde{\mathbf{P}}_i^{(j)} = \mathrm{exp}(\mathbf{S}_i^{(j)} - m_i)$; update $\ell_i = \mathrm{exp}(m_i^{\mathrm{old}} - m_i) \ell_i + \mathrm{rowsum}(\widetilde{\mathbf{P}}_i^{(j)})$. \label{code-ws:softmax_end}
\STATE Wait until $\mathbf{V}_j$ is loaded into shared memory.
\STATE Update output: $\mathbf{O}_i = \mathrm{diag}(\mathrm{exp}(m_i^{\mathrm{old}} - m_i))^{-1} \mathbf{O}_i + \widetilde{\mathbf{P}}_i^{(j)} \mathbf{V}_j$ utilizing RS-GEMM. Commit, then wait. \label{alg:ws_only_gemm2}
\STATE Mark the $(j\,\%\,s)$ buffer stage as released to the producer.
\ENDFOR
\STATE Finalize output: compute $\mathbf{O}_i = \mathrm{diag}(\ell_i)^{-1} \mathbf{O}_i$; set $L_i = m_i + \log(\ell_i)$.
\STATE Write the computed $\mathbf{O}_i$ and $L_i$ to HBM as the $i$th segment of outputs $\mathbf{O}$ and $L$.
\ENDIF
\end{algorithmic}
\end{algorithm}

In our deployment of \cref{alg:flash3_wgmma_ws_only} on the Hopper architecture, we employ \verb|setmaxnreg| to manage (de)allocations, utilize TMA to load $\mathbf{Q}_i$ and the set $\{ \mathbf{K}_j, \mathbf{V}_j \}_{0 \leq j < T_c}$, and leverage WGMMA for carrying out the GEMMs within the consumer mainloop. The SS and RS prefixes specify whether the primary operand resides in shared memory or the register file, respectively. When examining the execution sequence described in \cref{alg:flash3_wgmma_ws_only}, it should be observed that TMA load instructions proceed asynchronously, such that subsequent loads are not delayed by previous ones. Additionally, within the producer mainloop, the first $s$ iterations do not enforce waits, as they are primarily involved in filling the buffer.

\paragraph{Pingpong scheduling} The inherent asynchrony of WGMMA and TMA, combined with the use of warp specialization, enables the possibility of executing softmax computations for one warpgroup simultaneously with GEMM computations in another. This possibility is highlighted by the substantial disparity in throughput between matmul and non-matmul operations on contemporary hardware accelerators. For instance, the H100 SXM5 GPU can deliver 989 TFLOPS in FP16 matmul operations, whereas its peak throughput for special functions like exponentiation\footnote{According to the CUDA programming guide, each streaming multiprocessor (SM) is capable of 16 special function operations per clock cycle. Given 132 SMs and a clock frequency of 1830 MHz, this results in a theoretical throughput of 3.9 TFLOPS for special functions.}—which are essential for softmax—reaches only 3.9 TFLOPS. In the context of the FP16 attention forward pass with a 128-dimensional head, the volume of matmul FLOPS surpasses the requirement for exponentials by a factor of 512, yet the exponential computation is limited by 256-fold lower throughput, leading to scenarios where the exponential operations occupy up to 50\% of the total compute time when compared to matmul. This imbalance is further exacerbated with FP8 arithmetic, as matmul performance doubles while the rate of exponentials remains unchanged.

Because the exponential operation is executed by a dedicated hardware unit (the multi-function unit), optimal scheduling would entail performing exponentiation computations in parallel with matmuls executed by the Tensor Cores. To accomplish this, we insert synchronization barriers (\texttt{bar.sync} instructions) to ensure the GEMMs (specifically, GEMM1 corresponding to $\mathbf{P} \mathbf{V}$ in the current iteration and GEMM0 representing $\mathbf{Q} \mathbf{K}^\top$ in the upcoming iteration) within warpgroup 1 are dispatched ahead of those in warpgroup 2. This arrangement causes warpgroup 1's softmax to be processed concurrently with GEMMs being performed by warpgroup 2. Afterward, the two warpgroups switch tasks—warpgroup 2 computes softmax while warpgroup 1 executes GEMMs—thus establishing a “pingpong” scheduling pattern as shown in~\cref{fig:pingpong_scheduling}. While this pingpong execution does not strictly alternate as the figure may suggest, our experiments generally observe it yields notable performance gains (e.g., increasing from 570 TFLOPS to between 620 and 640 TFLOPS for FP16 forward passes with a head dimension of 128 and a sequence length of 8192).

\paragraph{Attention variants} For multi-query attention~\citep{shazeer2019fast} and grouped query attention~\citep{ainslie2023gqa}, our method aligns with that used in \textsc{FlashAttention-2}, modifying tensor indexing so as to prevent redundancy of $\mathbf{K}$ and $\mathbf{V}$ in HBM.

\subsection{Intra-warpgroup overlapping GEMMs and softmax}

It is possible to achieve partial instruction overlap between the softmax operations and GEMMs within a single warpgroup. Here, we outline a method for accomplishing this.

In the attention algorithm, operations within the inner loop (main loop) have sequential dependencies that impede parallelization within a single iteration. For example, (local) softmax (lines \ref{code-ws:softmax_start} to \ref{code-ws:softmax_end}) relies on the output $\mathbf{S}_i^{(j)}$ of the first GEMM, while the second GEMM takes its result $\widetilde{\mathbf{P}}_i^{(j)}$ as an operand. Indeed, the wait statements in lines \ref{alg:ws_only_gemm1} and \ref{alg:ws_only_gemm2} of \cref{alg:flash3_wgmma_ws_only} serialize the execution of softmax and GEMMs. However, we can break these dependencies by pipelining across iterations through additional buffers in registers. Pursuing this idea, we propose the following two-stage\footnote{Note that the number of stages of the overlapping scheme is bounded by, but need not equal, the number $s$ of stages in the circular SMEM buffer.} GEMM-softmax pipelining algorithm:

\begin{algorithm}[H]
    \caption{\small\label{alg:flash3_wgmma}\textsc{FlashAttention-3} consumer warpgroup forward pass}
    \begin{algorithmic}[1]
      \REQUIRE  Matrices $\mathbf{Q}_i \in \mathbb{R}^{B_r \times d}$ along with $\mathbf{K}, \mathbf{V} \in \mathbb{R}^{N \times d}$ stored in HBM, key block size $B_c$ given, where $T_c = \lceil \frac{N}{B_c} \rceil$.
      \STATE Adjust the number of registers in use according to the count of consumer warps involved.
      \STATE Allocate on-chip memory to $\mathbf{O}_i = (0) \in \mathbb{R}^{B_r \times d}$ and set $\ell_i, m_i = (0), (-\infty)$ with $\ell_i, m_i \in \mathbb{R}^{B_r}$.
      \STATE Stall computation until $\mathbf{Q}_i$ and the initial $\mathbf{K}_0$ have been transferred to shared memory.
      \STATE Perform WGMMA to evaluate $\mathbf{S}_{\mathrm{cur}} = \mathbf{Q}_i \mathbf{K}_0^T$, then commit the result and synchronize.
      \STATE Free the buffer’s first stage allocated to $\mathbf{K}$.
      \STATE Update $m_{i}$, compute $\tilde{\mathbf{P}}_{\mathrm{cur}}$, and adjust $\ell_{i}$ based on $\mathbf{S}_{\mathrm{cur}}$; rescale the values in $\mathbf{O}_i$ accordingly.
      \FOR{$1 \le j < T_c - 1$}
        \STATE Await arrival of $\mathbf{K}_j$ into shared memory.
        \label{code:mainloop_start}
        \STATE Utilize WGMMA to derive $\mathbf{S}_{\mathrm{next}} = \mathbf{Q}_i \mathbf{K}_{j}^T$; commit the result without blocking. \label{code:first_wgmma}
        \STATE Pause until $\mathbf{V}_{j-1}$ is available in shared memory.
        \STATE Calculate $\mathbf{O}_{i} = \mathbf{O}_{i} + \tilde{\mathbf{P}}_{\mathrm{cur}} \mathbf{V}_{j-1}$ using WGMMA, committing this update without blocking execution. \label{code:second_wgmma}
        \STATE Synchronize, ensuring WGMMA computation $\mathbf{Q}_i \mathbf{K}_{j}^T$ is finished.
        \STATE Based on $\mathbf{S}_{\mathrm{next}}$, update $m_{i}$, $\tilde{\mathbf{P}}_{\mathrm{next}}$, and $\ell_{i}$. \label{code:softmax}
        \STATE After confirming $\tilde{\mathbf{P}}_{\mathrm{cur}} \mathbf{V}_{j-1}$ (via WGMMA) is complete, carry out rescaling of $\mathbf{O}_i$.
        \STATE Release stage $(j\,\%\,s)$ of the $\mathbf{K}$ buffer and stage $((j-1)\,\%\,s)$ for $\mathbf{V}$.
        \STATE Transfer $\mathbf{S}_{\mathrm{next}}$ into $\mathbf{S}_{\mathrm{cur}}$ to proceed with the loop.
        \label{code:mainloop_end}
      \ENDFOR \STATE Stall until $\mathbf{V}_{T_c - 1}$ is loaded into shared memory.
      \STATE With WGMMA, evaluate and finalize
      $\mathbf{O}_{i} = \mathbf{O}_{i} + \tilde{\mathbf{P}}_{\mathrm{last}} \mathbf{V}_{T_c - 1}$, committing the operation and synchronizing upon completion.

\STATE Finalize: Adjust $\mathbf{O}_{i}$ according to $m_{i}$. Derive $L_{i}$ utilizing $m_{i}$ and $\ell_{i}$.
Store $\mathbf{O}_{i}$ and $L_{i}$ in HBM as the $i$-th entries of $\mathbf{O}$ and $L$ respectively.
\end{algorithmic}
\end{algorithm}

\cref{alg:flash3_wgmma} serves as the counterpart for the consumer pathway in \cref{alg:flash3_wgmma_ws_only}, together constituting the full implementation of the \textsc{FlashAttention-3} algorithm operating in FP16. Conceptually, we refer to WGMMA as synonymous with asynchronous GEMM. During execution of the main loop (spanning lines \ref{code:mainloop_start} through \ref{code:mainloop_end}), the algorithm concurrently performs the second WGMMA instruction of the current iteration $j$ (at line \ref{code:second_wgmma}) alongside the softmax computations associated with the subsequent iteration $j+1$ (at line \ref{code:softmax}).

Although the pipelined framework shown previously can, in principle, provide considerable performance improvements, several pragmatic issues arise:

\paragraph{Compiler reordering} While the pseudocode delineates an ideal pipeline execution, the actual instruction stream produced by the compiler (NVCC) may be subject to reordering for optimization purposes. Such reordering can disrupt the intended interleaving of WGMMA and non-WGMMA operations, possibly undermining the anticipated pipeline benefits and leading to unpredictable results. Examination of the corresponding SASS output reveals that the compiler generally does generate the expected interleaved code (see Section~\ref{sec:2-stage-sass}).

\paragraph{Register pressure} To achieve high throughput, minimizing register spills is crucial. However, the 2-stage pipelining approach inherently introduces greater demand on register resources, as intermediate results and pipeline context—most notably, an extra $\mathbf{S}_{\mathrm{next}}$—must be maintained in registers. Each threadblock consequently sees a rise in register allocation, specifically by $B_r \times B_c \times \text{sizeof}(\text{float})$. This heightened register utilization may restrict the feasible size of threadblocks, which is also a performance-relevant parameter and typically competes for registers. Practical deployment should therefore make decisions on tile size and pipeline depth informed by profiling and experimentation.

\paragraph{3-stage pipelining} Building on the two-stage scheme, a three-stage pipelined approach can further increase overlap between the second WGMMA and the softmax computation. While this has the theoretical advantage of boosting Tensor Core occupancy, it also increases register requirements, as it introduces one more pipeline stage that must be resident in registers at all times. Balancing the expanded register usage with optimal block or tile sizes thus becomes more complex. Further algorithmic details and empirical analysis of this approach appear in ~\cref{sec:3-stage}.

\subsection{Low-precision with FP8}
\label{sec:algofp8}

\textbf{Efficiency: layout transformations.} Executing the forward pass of \textsc{FlashAttention-3} using FP8 precision introduces extra complexities related to layout compatibility that are not present when using FP16.

To begin, it is important to highlight that input tensors $\mathbf{Q}$, $\mathbf{K}$, and $\mathbf{V}$ are commonly stored such that the head dimension is contiguous. However, for the second GEMM in FP8 WGMMA, the k-major layout necessitates that $\mathbf{V}$—specifically, the tiles of $\mathbf{V}$ that are transferred into SMEM—should be laid out contiguously along the sequence length dimension. As the TMA load operation is unable to modify which dimension is contiguous, we are left with two approaches: (1) perform a transpose of $\mathbf{V}$ in GMEM as a pre-processing procedure, or (2) conduct a tile-level transpose of $\mathbf{V}$ within the kernel subsequent to their placement in SMEM. For the first method, this can be accomplished in either of two ways: (1a) the transpose can be fused into the epilogue of a prior stage, such as immediately following the rotary embedding, or (1b) a separate pre-processing transpose kernel can be invoked\footnote{A well-optimized transpose kernel can nearly saturate the available memory bandwidth~\citep{colfax_cutlass_transpose_2024}.} to switch the strides between the head and sequence length axes. Nevertheless, integrating option (1a) is nontrivial within a generic library context, and adopting (1b) introduces excessive overhead under memory-bound scenarios like inference.

Rather than using alternative strategies, for FP8 \textsc{FlashAttention-3} we select approach (2). During the in-kernel transpose, we exploit the capabilities of the LDSM (\verb|ldmatrix|) and STSM (\verb|stmatrix|) instructions. These instructions enable a warp of threads to collaboratively load from shared memory (SMEM) to registers (RMEM) and to store from RMEM back to SMEM, both at 128-byte blocks. \footnote{The PTX manual specifies LDSM and STSM as operating on $8 \times 8$ matrices with 16-bit elements \cite[\S 9.7.13.4.15-16]{ptx}, but by packing 8-bit values in pairs, we are able to leverage LDSM/STSM for FP8 data. Notably, since the transpose variants of LDSM/STSM cannot separate packed 8-bit values, some register shuffling between LDSM and STSM is necessary to carry out a true tile-wise transpose; the specific operations are not detailed here.} LDSM and STSM not only minimize register consumption—making it practical to use them in the producer warpgroup—but also facilitate transposes during memory transfer. In addition, subsequent to the initial iteration, the transposition of the following $\mathbf{V}$ tile can be overlapped with the two WGMMAs acting on the previous $\mathbf{V}$ tile and the current $\mathbf{K}$ tile, utilizing this computation window.

Second, we note that the memory organization of the FP32 accumulator in an FP8 WGMMA does not align with the arrangement anticipated for operand A when stored in registers, in contrast to the scenario with FP16. Illustrations of partial layouts for these two cases are provided in \cref{fig:rmem_accum} and \cref{fig:rmem_operand}, showing how register entries are assigned per thread according to the ordering given. By leveraging byte permutation instructions, the accumulator output from the initial WGMMA can be reformatted to meet the input requirements of a subsequent WGMMA and to correspond with the organization of the $\mathbf{V}$ tile generated by the in-kernel transpose operation. Concretely, as depicted in \cref{fig:rmem_accum}, we modify the sequence to
$$\{ \verb|d0 d1 d4 d5 d2 d3 d6 d7| \},$$
and this register sequence repeats for every group of 8 bytes. At the logical level, this operation effectively permutes the columns of the $\mathbf{P}$ tile (for instance, columns $0189$ are mapped to occupy the initial four columns). To ensure WGMMA produces the intended output tile, we may organize the in-kernel transpose to emit the $\mathbf{V}$ tile with a corresponding row reordering.\footnote{This extra flexibility in the in-kernel transpose makes it unnecessary to employ shuffle instructions for reassigning registers across threads, a step we had previously discussed in~\citep{colfax_fp8_flashattention_2024}.}

\textbf{Accuracy: block quantization and incoherent processing.} The FP8 (e4m3) representation allocates just 3 bits for the mantissa and 4 bits for the exponent, which introduces greater numerical inaccuracies compared to FP16 or BF16 formats. Additionally, large-scale models tend to exhibit outlier values~\citep{dettmers2208llm, sun2024massive} with magnitudes significantly exceeding those of most elements, further complicating the quantization procedure. Standard practice involves applying per-tensor scaling~\citep{micikevicius2022fp8}, where a single scaling factor is maintained for each tensor (e.g., separately for $\mathbf{Q}$, $\mathbf{K}$, and $\mathbf{V}$). To mitigate the numerical errors introduced by attention computations in FP8, we adopt two strategies:

\begin{enumerate}

\item \textbf{Block quantization}: In this approach, a single scalar is maintained for each block. Specifically, for tensors $\mathbf{Q}$, $\mathbf{K}$, and $\mathbf{V}$, we divide each into separate blocks of dimensions $B_r \times d$ or $B_c \times d$, and quantize each block individually. This quantization step can be integrated with pre-attention operations, such as rotary embedding, without incurring extra latency (because the process is limited by memory bandwidth). Since the \textsc{FlashAttention-3} algorithm already works on the block level, scaling can be applied to every block of $\mathbf{S}$ to reflect the quantization, introducing no additional computational burden.

\item \textbf{Incoherent processing}: To mitigate the impact of outliers, we pre-multiply $\mathbf{Q}$ and $\mathbf{K}$ by a random orthogonal matrix $\mathbf{M}$ prior to their FP8 quantization. Owing to the orthogonality property ($\mathbf{M} \mathbf{M}^\top = I$), the product $(\mathbf{Q} \mathbf{M}) (\mathbf{K} \mathbf{M})^\top$ is identical to $\mathbf{Q} \mathbf{K}^\top$, guaranteeing that the transformation does not alter the result of the attention mechanism. This procedure redistributes the influence of outlier values, as each component in $\mathbf{Q} \mathbf{M}$ or $\mathbf{K} \mathbf{M}$ comprises a random combination of components from $\mathbf{Q}$ or $\mathbf{K}$, leading to decreased quantization noise. Following the methods of \citet{chee2024quip} and~\citet{tseng2024quip}, we set $\mathbf{M}$ as a product of random diagonal matrices with entries $\pm 1$ and a Hadamard matrix, allowing for efficient multiplication in $O(d \log d)$ time rather than $O(d^2)$ and making it feasible to fuse this step with rotary embedding without incurring further computational overhead.
\end{enumerate}

Our experiments, presented in \cref{sec:numerical_error}, demonstrate that combining these strategies can lower numerical error by up to a factor of 2.6.

\section{Empirical Validation}
\label{sec:exp}

To implement \textsc{FlashAttention-3}, we leverage WGMMA and TMA abstractions provided by CUTLASS~\citep{Thakkar_CUTLASS_2023}, enabling us to assess both its performance and precision.

\begin{itemize}

\item \textbf{Benchmarking attention.} We evaluate the execution time of \textsc{FlashAttention-3} over a range of sequence lengths and benchmark its performance against the baseline PyTorch implementation, \textsc{FlashAttention-2}, the Triton variant of \textsc{FlashAttention-2} (which utilizes H100-specific kernel instructions), and a cuDNN-based, vendor-optimized version of \textsc{FlashAttention-2} for H100 GPUs. Our findings show that \textsc{FlashAttention-3} achieves speedups of up to 2.0$\times$ compared to \textsc{FlashAttention-2}, and up to 1.5$\times$ relative to the Triton version. Furthermore, \textsc{FlashAttention-3} attains a peak throughput of 740 TFLOPs/s, corresponding to 75\% of H100 GPUs’ peak theoretical TFLOPs/s.

\item \textbf{Ablation study.} We empirically demonstrate that the enhancements from warp-specialization and the GEMM-softmax pipeline are critical contributors to the acceleration observed with \textsc{FlashAttention-3}.

\item \textbf{Accuracy of FP8 attention.} Through experimental validation, we establish that employing block quantization together with incoherent processing effectively reduces FP8 \textsc{FlashAttention-3}’s numerical errors by a factor of 2.6$\times$.
\end{itemize}

\subsection{Benchmarking Attention}
\label{subsec:benchmark_attn}

We evaluate the execution times of several attention mechanisms on an H100 80GB SXM5 GPU across varying configurations, including presence or absence of a causal mask and head sizes of 64 or 128, utilizing FP16 input tensors. The benchmarks, presented in~\cref{fig:benchmark_attn_fwd} and~\cref{fig:benchmark_attn_bwd}, indicate that \textsc{FlashAttention-3} achieves roughly 1.5 to 2 times faster forward computations and 1.5 to 1.75 times faster backward computations relative to \textsc{FlashAttention-2}. Relative to the conventional attention approach, \textsc{FlashAttention-3} delivers speedups in the range of 3 to 16 times. Notably, for moderate to extended sequence lengths (1k tokens or more), the performance of \textsc{FlashAttention-3} even exceeds that of a closed-source vendor library (cuDNN) specifically optimized for the H100 GPUs.

\paragraph{Benchmark settings:} We vary the sequence length as 512, 1k, ..., 16k, and set batch size so that the total number of tokens is 16k. We set the hidden dimension to 2048, and head dimension to be either 64, 128, or 256 (i.e., 32 heads, 16 heads, or 8 heads). To calculate the FLOPs of the forward pass, we use:

\begin{equation*}
  4 \cdot \text{seqlen}^2 \cdot \text{head dimension} \cdot \text{number of heads}.
\end{equation*}

Applying causal masking reduces this figure by half, as roughly 50% of the computations are performed. For the backward pass, we estimate its FLOPs by multiplying those of the forward pass by 2.5, reflecting the ratio of 5 matrix multiplications during backward propagation (with recomputation) to the 2 in the forward phase.

Additionally, we assess the runtime of FP8 during the forward pass using comparable experimental conditions. The outcomes for a headdim value of 256 are presented in ~\cref{fig:benchmark_attn_fp8}, while comprehensive results can be found in \cref{sec:benchmark_attn_fp8_full}.

\subsection{Ablation Study: 2-Stage Pipelining Experiments}
\label{sec:2-stage-pipelining-experiments}

We conduct an ablation study on both the two-stage WGMMA-softmax pipelining and warp-specialization within non-causal FP16 \textsc{FlashAttention-3}, maintaining fixed parameters $\{\text{batch}, \text{seqlen}, \text{nheads}, \text{hdim}\} = \{ 4, 8448, 16, 128\}$. The findings presented in~\cref{table:ablation_pipelining} demonstrate that our proposed algorithmic strategies—incorporating warp-specialization and enabling asynchrony by overlapping GEMM computation with the softmax operation—deliver substantial performance gains, improving throughput from 570 TFLOPs to 661 TFLOPs.

\subsection{Numerical Error Validation}
\label{sec:numerical_error}

Given growing attention to the numerical inaccuracies in \textsc{FlashAttention}~\citep{golden2024flash}, we evaluate \textsc{FlashAttention-2}, \textsc{FlashAttention-3}, and a baseline attention implementation by comparing them with a reference computed in FP64. In order to mimic outlier activations and features observed in LLMs~\citep{dettmers2208llm, sun2024massive}, we sample the elements of $\mathbf{Q}, \mathbf{K}, \mathbf{V}$ according to the following distribution:

\begin{equation*}
  \mathcal{N}(0, 1) + \mathcal{N}(0, 100) \cdot \mathrm{Bernoulli}(0.001).
\end{equation*}

Specifically, the entries are drawn from a normal distribution with zero mean and unit standard deviation, except that for 0.1\% of them, an independent Gaussian noise with standard deviation 10 is added. The root mean squared error (RMSE) is subsequently assessed, as shown in~\cref{table:numerical_error}. When operating in FP16, both \textsc{FlashAttention-2} and \textsc{FlashAttention-3} reduce RMSE by a factor of 1.7 compared to the conventional approach, attributable to storing intermediate (softmax) results in FP32 precision. For the FP8 baseline, the attention mechanism employs per-tensor scaling; its matmul accumulator works in FP32 and the intermediate softmax values are stored in FP16. With block quantization and incoherent processing strategies, \textsc{FlashAttention-3} in FP8 demonstrates 2.6$\times$ improved accuracy over this baseline method.

\section{Dicussion, Limitations, Conclusion}
\label{sec:discussion}

Through \textsc{FlashAttention-3}, we show that advances in programming paradigms and hardware capabilities—such as asynchronous execution and reduced precision computation—can substantially enhance both the speed and fidelity of attention mechanisms. Our approach achieves a 1.5-2.0$\times$ increase in attention throughput over \textsc{FlashAttention-2} and decreases FP8 numerical errors by a factor of 2.6 compared with traditional per-tensor quantization methods. There remain, however, several areas for improvement that we intend to pursue. These include further optimization for LLM inference workloads, implementing a persistent kernel architecture for the FP8 version,\footnote{In our benchmarks, only the FP16 variant of \textsc{FlashAttention-3} employs a persistent kernel and advanced load balancing, whereas the FP8 implementation lacks these optimizations. This is a key reason why FP8 \textsc{FlashAttention-3} underperforms in scenarios with short sequences and causal masks, relative to FP8 cuDNN kernels.} and deeper investigation into the behavior of low-precision attention during large-scale model training. Although this work centers on Hopper GPUs, we anticipate that our techniques will extend effectively to a range of other accelerators. Ultimately, we hope that providing a faster and more precise attention primitive will foster new opportunities for tasks requiring extended context.

\end{document}